\chapter*{Abstract}
\begin{otherlanguage}{ngerman}
\section*{Deutsch}
Interpretierbarkeit entwickelt sich von der nachträglichen Analyse einzelner Vorhersagen zu einem Werkzeug, mit dem Modelle gezielt verbessert und gesteuert werden. Wenn Attributionen die Korrektur unerwünschter Merkmale, das Training oder die Inferenz beeinflussen, muss deshalb auch ihre Berechnung verstanden und überprüft werden. Diese Arbeit verbindet dafür zwei Perspektiven: die strukturierte Darstellung interner Beiträge und die mathematische Beschreibung ihrer Rückwärtspropagation.

Ein trainiertes neuronales Netz zu verstehen hei\ss{}t, nachzuvollziehen, wie es eine Eingabe in eine Ausgabe \"uberf\"uhrt --- welche gegenl\"aufigen Einfl\"usse jede Vorhersage gegeneinander abw\"agt. Diese Arbeit untersucht Interpretierbarkeit durch \emph{Zerlegung}: Statt ein Netz zur Interpretierbarkeit zu zwingen, legen wir interpretierbare Struktur offen, die ein trainiertes ReLU-Netz bereits enth\"alt. Ein einziger Begriff der \emph{Netzzerlegung} --- eine Anhebung der Berechnung in ein Paar koordinierter, nichtnegativer Str\"ome, deren Differenz das urspr\"ungliche Netz exakt reproduziert --- vereinheitlicht mehrere solcher Strukturen, ohne eine einzige Vorhersage zu ver\"andern.

Die erste ist die Differenz-konvexer-Funktionen-Zerlegung (DC-Zerlegung). Da Monotonie die Interpretierbarkeit f\"ordert, aber die Ausdrucksst\"arke einschr\"ankt, trainieren wir kein monotones Modell, sondern zerlegen ein trainiertes ReLU-Netz --- ein MLP oder ein CNN --- in zwei Teilnetze $g$ und $h$ mit $f = g - h$, die in den verdoppelten Eingabekoordinaten konvex und monoton sind und positive und negative Evidenz auf getrennten Pfaden f\"uhren. Wir stabilisieren die ausgewertete ReLU-Zerlegung und konstruieren aus der absoluten Kr\"ummung exakte Paare f\"ur glatte Aktivierungen (GELU, Swish, Tanh, Sigmoid, Softplus). Explizite Softmax-, Produkt- und LayerNorm-Paare mit positivem $\varepsilon$ ergeben zudem eine DC-Zerlegung deterministischer Transformer beliebiger endlicher Tiefe; daraus folgt keine Garantie numerisch stabiler Auswertung. Auf der ReLU-Zerlegung bauen wir die Saliency-Methoden SplitCAM und SplitLRP auf, die auf VGG-16 und ResNet-18 \"uber ImageNet-S konfigurationsabh\"angige Verbesserungen bei Lokalisierung und Selectivity sowie Zielkonflikte bei Pixel Flipping und Robustheit zeigen. Zudem trainieren wir Differenzen konvexer oder monotoner Netze zur MNIST-Klassifikation und beobachten klassenabh\"angige Muster in ihren Stromgradienten als qualitative Illustration der Darstellung.

Die zweite ist die Stopp-Zerlegung, die die Gewichte nach Vorzeichen aufteilt, $W = W^+ - W^-$, und exzitatorische von inhibitorischen Beitr\"agen entlang der Netzpfade trennt. Diese Zerlegung ist die Br\"ucke zu einer spieltheoretischen Deutung des R\"uckw\"artsdurchlaufs. Aufbauend auf dem ReLU-Netz-Spiel von Gaubert und Vlassopoulos fassen wir R\"uckw\"arts-Attribution als ein Zwei-Spieler-Nullsummen-Stoppspiel auf einem erweiterten Netzgraphen auf. Gradienten und die vollst\"andige $\alpha\beta$-LRP-Familie ergeben sich als Integrale \"uber Spieltrajektorien unter bestimmten Gleichgewichten, sodass Attributionskarten zu Projektionen einer Verteilung \"uber Trajektorien werden, statt das prim\"are Objekt zu sein. Erw\"unschte Eigenschaften einer Interpretation --- Lokalisierungssch\"arfe, Robustheit gegen\"uber Eingangsrauschen, stabiles Attention-Routing --- lassen sich als spieltheoretische Begriffe (Politik-Regularisierung, Risikoaversion, erweiterte Aktionsr\"aume) formulieren, die sich unmittelbar in neue R\"uckw\"artsregeln \"ubersetzen; auf ViT-B/16 \"ubertrifft eine solche Anpassung von $\alpha\beta$-LRP alle betrachteten Vergleichsmethoden, einschlie\ss{}lich transformer-spezifischer Methoden, bei der Attributionslokalisierung. Die revidierte Evaluation trennt Lokalisierung und Faithfulness bei der Validierungsauswahl und berichtet Konfidenzintervalle sowie gepaarte Tests; andere Metriken zeigen Zielkonflikte.

Ein einziger Rahmen macht so mehrere Lesarten derselben trainierten Gewichte zugleich verf\"ugbar. Wir entwickeln zudem die tropisch-geometrische Perspektive, in der die zerlegten Str\"ome tropische Polynome sind, deren Newton-Polytope --- \emph{imaginative Polytope} --- die erreichbaren Gradienten als gelernte visuelle Prototypen kodieren, sowie eine Trajektorienverteilungs-Sicht, die den R\"uckw\"artsdurchlauf mit einer Hellinger-Distanz zwischen normalisierten Routing-Verteilungen ausstattet und damit das Routing der Trajektorien vor der modellspezifischen Projektion signierter, diskontierter Belegungsmaße auf Heatmaps vergleicht.
\\[2ex]
\textbf{Schl\"usselw\"orter:} \itshape \germankeywords
\end{otherlanguage}

\begin{otherlanguage}{english}
\section*{English}
Interpretability is moving from post hoc analysis of individual predictions toward improving and controlling models. When attributions guide the correction of spurious features, training, or inference, the calculation producing them must itself be understood and checked. This thesis connects two perspectives on that problem: a structured representation of internal contributions and a mathematical account of their backward propagation.

Understanding a trained neural network means accounting for how it turns an input into an output --- for the competing influences that every prediction balances against one another. This thesis studies interpretability through \emph{decomposition}: instead of constraining a network to make it interpretable, we expose interpretable structure that a trained ReLU network already contains. A single notion of a \emph{network decomposition} --- a lift of the computation into a pair of coordinated, non-negative streams whose difference reproduces the original network exactly --- unifies several such structures, none of which alters a single prediction.

The first is the difference-of-convex (DC) decomposition. Monotonicity is known to aid interpretability but to restrict expressivity, so rather than training a monotone model we split a trained ReLU network --- an MLP or a CNN --- into two subnetworks $g$ and $h$ with $f = g - h$, convex and monotone in doubled input coordinates and carrying positive and negative evidence on separate pathways. We stabilize the evaluated ReLU split and use absolute curvature to construct exact pairs for smooth activations (GELU, Swish, tanh, sigmoid, softplus). Explicit softmax, product, and positive-epsilon LayerNorm pairs also give a DC decomposition of deterministic Transformers of arbitrary finite depth, without a numerical-stability guarantee. On the ReLU split representation we build the saliency methods SplitCAM and SplitLRP, which show configuration-dependent gains in localization and Selectivity on VGG-16 and ResNet-18 over ImageNet-S, with trade-offs in Pixel Flipping and robustness. We also train differences of convex or monotone networks for MNIST classification and observe class-dependent patterns in their stream gradients, providing a qualitative illustration of the representation.

The second is the stopping decomposition, which splits the weights by sign, $W = W^+ - W^-$, separating excitatory from inhibitory contributions along the paths of the network. This decomposition is the bridge to a game-theoretic reading of the backward pass. Building on Gaubert and Vlassopoulos' ReLU Net Game, we recast backward attribution as a two-player zero-sum stopping game on an extended network graph. Gradients and the full $\alpha\beta$-LRP family arise as integrals over game trajectories under specific equilibria, so that attribution maps become projections of a distribution over trajectories rather than the primary object. Desired properties of an interpretation --- localization focus, robustness to input noise, stable attention routing --- can be stated as game-theoretic concepts (policy regularization, risk aversion, extended action sets) that translate directly into new backward rules; on ViT-B/16, one such adaptation of $\alpha\beta$-LRP outperforms all evaluated baselines, including transformer-specific methods, in attribution localization. The revised evaluation selects localization and faithfulness configurations independently on validation data and reports confidence intervals and paired tests; other metrics expose trade-offs.

One framework thus makes several readings of the same trained weights available at once. We develop the tropical-geometric perspective, in which the split streams are tropical polynomials whose Newton polytopes --- \emph{imaginative polytopes} --- encode the attainable gradients as learned visual prototypes; and a trajectory-distribution view that equips the backward pass with a Hellinger distance between normalized routing laws, comparing trajectory routing before the model-dependent projection of signed discounted occupation measures onto heatmaps.
\\[2ex]
\textbf{Keywords:} \itshape \englishkeywords
\end{otherlanguage}
